% Documento generado automáticamente desde notebooks/fase1.ipynb.
% Pensado para integrarse en Plantilla_Latex_GCD/tfgs/tex/resultados.tex.
\section{Resultados de la Fase 1}

\subsection{1. Objetivo de la fase}
Evaluar estructura, calidad, señal estadística, redundancia, dimensionalidad y riesgos de los datasets crudos antes de aplicar preprocesado, selección de características o modelado.
\subsection{2. Datasets analizados}
\begin{itemize}
\item \texttt{breast\_cancer\_wisconsin}: 569 filas, 32 columnas, target \texttt{target}.
\item \texttt{customer\_churn}: 440832 filas, 12 columnas, target \texttt{Churn}.
\item \texttt{madelon}: 2000 filas, 501 columnas, target \texttt{target}.
\item \texttt{olive\_oil}: 572 filas, 12 columnas, target \texttt{target}.
\end{itemize}

\subsection{3. Resumen estructural}
\begin{itemize}
\item \texttt{breast\_cancer\_wisconsin}: 30 features analíticas, 1 posible(s) identificador(es), ratio features/muestras 0.05272.
\item \texttt{customer\_churn}: 10 features analíticas, 1 posible(s) identificador(es), ratio features/muestras 2.268e-05.
\item \texttt{madelon}: 500 features analíticas, 0 posible(s) identificador(es), ratio features/muestras 0.25.
\item \texttt{olive\_oil}: 10 features analíticas, 1 posible(s) identificador(es), ratio features/muestras 0.01748.
\end{itemize}

\subsection{4. Calidad del dato}
\begin{itemize}
\item \texttt{breast\_cancer\_wisconsin}: 0 columnas con nulos, 0 duplicados, 0 constantes y 0 variables de baja variabilidad/dominancia.
\item \texttt{customer\_churn}: 0 columnas con nulos, 0 duplicados, 0 constantes y 9 variables de baja variabilidad/dominancia.
\item \texttt{madelon}: 0 columnas con nulos, 0 duplicados, 0 constantes y 24 variables de baja variabilidad/dominancia.
\item \texttt{olive\_oil}: 0 columnas con nulos, 0 duplicados, 0 constantes y 1 variables de baja variabilidad/dominancia.
\end{itemize}

\subsection{5. Distribución del target}
\begin{itemize}
\item \texttt{breast\_cancer\_wisconsin}: 2 clases; clase minoritaria n=212; ratio 1.68.
\item \texttt{customer\_churn}: 2 clases; clase minoritaria n=190833; ratio 1.31.
\item \texttt{madelon}: 2 clases; clase minoritaria n=1000; ratio 1.
\item \texttt{olive\_oil}: 9 clases; clase minoritaria n=25; ratio 8.24.
\end{itemize}

\subsection{6. Señal variable-target}
\begin{itemize}
\item \texttt{breast\_cancer\_wisconsin}: 27 variables candidatas univariantes tras FDR y efecto mínimo.
\item \texttt{customer\_churn}: 9 variables candidatas univariantes tras FDR y efecto mínimo.
\item \texttt{madelon}: 13 variables candidatas univariantes tras FDR y efecto mínimo.
\item \texttt{olive\_oil}: 10 variables candidatas univariantes tras FDR y efecto mínimo.
\end{itemize}

\subsection{7. Corrección por múltiples comparaciones}
\begin{itemize}
\item \texttt{breast\_cancer\_wisconsin}: 27 significativas sin corregir frente a 27 tras FDR.
\item \texttt{customer\_churn}: 10 significativas sin corregir frente a 10 tras FDR.
\item \texttt{madelon}: 38 significativas sin corregir frente a 13 tras FDR.
\item \texttt{olive\_oil}: 10 significativas sin corregir frente a 10 tras FDR.
\end{itemize}

\subsection{8. Tamaños de efecto}
\begin{itemize}
\item \texttt{breast\_cancer\_wisconsin}: grande=25, medio=1, muy\_pequeño=2, pequeño=2.
\item \texttt{customer\_churn}: grande=4, medio=3, muy\_pequeño=1, pequeño=2.
\item \texttt{madelon}: medio=7, muy\_pequeño=459, pequeño=34.
\item \texttt{olive\_oil}: grande=10.
\end{itemize}

\subsection{9. Redundancia}
\begin{itemize}
\item \texttt{breast\_cancer\_wisconsin}: 29 pares con |Spearman| >= 0.85.
\item \texttt{customer\_churn}: 0 pares con |Spearman| >= 0.85.
\item \texttt{madelon}: 12 pares con |Spearman| >= 0.85.
\item \texttt{olive\_oil}: 2 pares con |Spearman| >= 0.85.
\end{itemize}

\subsection{10. Dimensionalidad y ruido}
\begin{itemize}
\item \texttt{breast\_cancer\_wisconsin}: riesgo bajo; posible ruido univariante 4 variables.
\item \texttt{customer\_churn}: riesgo bajo; posible ruido univariante 3 variables.
\item \texttt{madelon}: riesgo medio; posible ruido univariante 493 variables.
\item \texttt{olive\_oil}: riesgo bajo; posible ruido univariante 0 variables.
\end{itemize}

\subsection{11. Posibles relaciones espurias o leakage}
\begin{itemize}
\item \texttt{breast\_cancer\_wisconsin}: 1 variables para revisión semántica; no se confirma leakage en esta fase.
\item \texttt{customer\_churn}: 1 variables para revisión semántica; no se confirma leakage en esta fase.
\item \texttt{madelon}: 25 variables para revisión semántica; no se confirma leakage en esta fase.
\item \texttt{olive\_oil}: 2 variables para revisión semántica; no se confirma leakage en esta fase.
\end{itemize}

\subsection{12. Preclasificación preliminar}
\begin{itemize}
\item \texttt{breast\_cancer\_wisconsin}: candidata\_fuerte=25, candidata\_moderada=1, pendiente\_multivariante=1, posible\_ruido\_univariante=2, sospechosa=1.
\item \texttt{customer\_churn}: candidata\_fuerte=7, candidata\_moderada=2, sospechosa=1.
\item \texttt{madelon}: candidata\_fuerte=7, pendiente\_multivariante=3, posible\_ruido\_univariante=453, redundante\_o\_correlacionada=12, sospechosa=25.
\item \texttt{olive\_oil}: candidata\_fuerte=8, sospechosa=2.
\end{itemize}

\subsection{13. Figuras más importantes}
\begin{itemize}
\item \texttt{results/figures/01\_raw\_eda/01\_02\_estructura/raw\_structure\_summary.png}: Escala muestral, dimensionalidad inicial y ratio variables/muestras.
\item \texttt{results/figures/01\_raw\_eda/01\_03\_calidad/raw\_quality\_issues\_summary.png}: Incidencias accionables de calidad del dato crudo sin heatmap vacío.
\item \texttt{results/figures/01\_raw\_eda/01\_04\_target/raw\_target\_distribution\_panel.png}: Distribución de clases, con lectura separada para binarios y olive\_oil multiclase.
\item \texttt{results/figures/01\_raw\_eda/01\_09\_fdr/raw\_fdr\_before\_after.png}: Efecto de la corrección por múltiples comparaciones, especialmente en madelon.
\item \texttt{results/figures/01\_raw\_eda/01\_11\_redundancia/raw\_high\_corr\_pairs\_by\_dataset.png}: Redundancia potencial por correlación alta y prioridad de control en Fase 2.
\item \texttt{results/figures/01\_raw\_eda/01\_15\_ead/raw\_dataset\_profile\_heatmap.png}: Matriz de severidad EAD trazada desde raw\_dataset\_profiles.csv.
\end{itemize}

\subsection{14. Tablas más importantes}
\texttt{raw\_structure\_summary.csv}, \texttt{raw\_quality\_summary.csv}, \texttt{raw\_target\_distribution.csv}, \texttt{raw\_feature\_target\_tests.csv}, \texttt{raw\_fdr\_corrected\_tests.csv}, \texttt{raw\_effect\_sizes.csv}, \texttt{raw\_high\_correlation\_pairs.csv}, \texttt{raw\_feature\_preclassification.csv}.
\subsection{15. Conclusiones por dataset}
\begin{itemize}
\item \texttt{breast\_cancer\_wisconsin}: pasa a fases posteriores con riesgo dimensional bajo, 26 señales fuertes exploratorias y 1 riesgos a revisar.
\item \texttt{customer\_churn}: pasa a fases posteriores con riesgo dimensional bajo, 7 señales fuertes exploratorias y 1 riesgos a revisar.
\item \texttt{madelon}: pasa a fases posteriores con riesgo dimensional medio, 7 señales fuertes exploratorias y 25 riesgos a revisar.
\item \texttt{olive\_oil}: pasa a fases posteriores con riesgo dimensional bajo, 10 señales fuertes exploratorias y 2 riesgos a revisar.
\end{itemize}

\subsection{15.1 Matices de interpretación}
\begin{itemize}
\item \texttt{customer\_churn}: por su tamaño muestral, la significación estadística debe leerse junto al tamaño de efecto; no se interpreta como evidencia fuerte por p-value aislado.
\item \texttt{madelon}: el patrón sugiere alta dimensionalidad y mucho ruido univariante, pero no ausencia total de señal; hay variables con señal exploratoria tras FDR.
\item \texttt{olive\_oil}: queda marcado como problema multiclase con desbalance relevante; la evaluación posterior debe priorizar métricas macro o balanceadas.
\item \texttt{breast\_cancer\_wisconsin}: muestra señal univariante fuerte, pero también redundancia alta; la Fase 2 debe controlar variables correlacionadas antes de selección formal.
\end{itemize}

\subsection{16. Implicaciones para la Fase 2}
La Fase 2 debe usar explícitamente los artefactos de esta fase para tomar, como mínimo, las siguientes decisiones:
\begin{itemize}
\item IDs: excluir o aislar columnas con rol \texttt{posible\_id} antes de entrenar o seleccionar variables.
\item Baja varianza/dominancia: revisar \texttt{raw\_low\_variance\_features.csv} y decidir eliminación, recodificación o conservación justificada.
\item Codificación: codificar variables categóricas sin introducir información del target ni mezclar ajuste entre train y test.
\item Escalado: ajustar escaladores solo con train, especialmente para PCA, métodos basados en distancia y modelos sensibles a escala.
\item Outliers: revisar \texttt{raw\_outlier\_summary.csv} y decidir tratamiento robusto sin eliminar observaciones de forma automática.
\item Splits: mantener particiones estratificadas cuando proceda, con especial cuidado en \texttt{olive\_oil} por su target multiclase desbalanceado.
\item Leakage/proxies: auditar semánticamente \texttt{raw\_spurious\_risk\_features.csv} y \texttt{raw\_proxy\_leakage\_candidates.csv} antes de modelar.
\item Redundancia: usar \texttt{raw\_high\_correlation\_pairs.csv} para controlar bloques correlacionados, especialmente en \texttt{breast\_cancer\_wisconsin}.
\end{itemize}

\subsection{17. Incidencias y decisiones metodológicas}
No se inventan datos ni targets. Se excluyen identificadores de los análisis de asociación y correlación. Las visualizaciones vacías se omiten, especialmente mapas de nulos cuando no hay nulos.